\documentclass[../../../include/open-logic-section]{subfiles}

\begin{document}

\olfileid{sth}{ord-arithmetic}{expo}
\olsection{序数幂}

现在我们来讨论序数幂。遗憾的是，对于序数幂并不存在一个\emph{好的}构造性定义。

当然，\emph{确实有}显式的构造性定义。这里给出一种。令 $\text{finfun}(\alpha,\beta)$ 为所有满足 $\Setabs{\gamma \in
\alpha}{f(\gamma) \neq 0}$ 与某个自然数等势的函数 $f
\colon \alpha \to \beta$ 的集合。在 $\text{finfun}(\alpha,\beta)$ 上定义良序如下：$f
\sqsubset g$ 当且仅当 $f \neq g$ 且 $f(\gamma_0) < g(\gamma_0)$，其中
$\gamma_0 = \text{max}\Setabs{\gamma \in \alpha}{f(\gamma) \neq
g(\gamma)}$。于是我们可以将 $\ordexpo{\alpha}{\beta}$ 定义为
$\ordtype{\text{finfun}(\alpha, \beta), \sqsubset}$。Potter 采用了
这个显式定义，并立即解释道：

\begin{quote}
	选择该序关系纯粹是为了得到一个满足相应递归条件的序数幂定义……，并且它比序数和或序数积要难以想象得多。
	\citep[p.~199]{Potter2004}
\end{quote}

确实如此。我们曾将加法刻画为“一个 $\alpha$ 的副本后面跟着一个 $\beta$ 的副本”，将乘法刻画为“一个由 $\alpha$ 的副本组成的 $\beta$ 序列”。但对于 $\text{finfun}(\alpha, \gamma)$，我们则没有什么简洁的直观刻画。因此，我们将仅\emph{借由}超限递归来给出序数幂的定义，即：

\begin{defn}\ollabel{ordexporecursion}
\begin{align*}
	\ordexpo{\alpha}{0} &= 1\\
	\ordexpo{\alpha}{\beta\ordplus 1} &=\ordexpo{\alpha}{\beta} \ordtimes \alpha\\
	\ordexpo{\alpha}{\beta} &= \bigcup_{\delta < \beta}\ordexpo{\alpha}{\delta}& & \text{当 $\beta$ 是极限序数时}
\end{align*}
\end{defn}

如果我们纯粹\emph{以}集合论者的身份开展工作，或许会想探索序数幂的更多性质。但我们没有太多要补充的，只需指出一个并不令人意外的事实：序数幂不满足交换律。因此 $\ordexpo{2}{\omega} =
\bigcup_{\delta < \omega}\ordexpo{2}{\delta} = \omega$，而
$\ordexpo{\omega}{2} = \omega \ordtimes \omega$。不过话说回来，我们本就不该\emph{期望}幂运算满足交换律，因为自然数上的幂运算也不可交换：$\ordexpo{2}{3} = 8 < 9 = \ordexpo{3}{2}$。

\begin{prob}
使用超限归纳法证明：如果我们定义
$\ordexpo{\alpha}{\beta} = \ordtype{\text{finfun}(\alpha, \beta),
\sqsubset}$，则可以得到 \olref[sth][ord-arithmetic][expo]{ordexporecursion} 中的递归等式。
\end{prob}

\end{document}